\scanpage{168}
Si \uline{toutes} les données et tous les axiomes sont
réunies sur $\underline U$, on a ce qui suit ($k$ un \unsure{corps comm.}).
\[
\left\{\vcenter{\hbox{\begin{minipage}{0.88\linewidth}
$\underline U$ algèbre associative topologique complète, avec suffisamment de représentations continues de dim. finie

$\underline U$ muni d'une \deleted{\unsure{homom.|structure} d'algèbre \unsure{cont} $\underline U\xrightarrow{\Delta}\underline U\otimes\underline U$}\inserted{structure de co-algèbre associative} \deleted{\unsure{iden}}\inserted{et} commutative

Les deux structures compatibles\,: \unsure{c'est|$\Delta$ est} un hom. continu d'algèbre associative

$\underline U$ muni d'une augmentation $\varepsilon$ qui est un hom continu d'algèbre et un hom. de Co-Algèbre

$\underline U$ muni d'une involution $\vee$, qui est un antiautomorphisme \inserted{continu} d'algèbre associative, et un automorphisme de coalgèbre;

$\varepsilon$ \unsure{invariante|invariant} sous $\vee$\,: $\varepsilon(\breve x)=\varepsilon(x)$

On a \quad $P(\vee\otimes1)\Delta=P(1\otimes\vee)\Delta=1$
\end{minipage}}}\right.
\]
($P$ désigne en principe l'hom $\underline U\mathbin{\widehat\otimes}\underline U\longrightarrow\underline U$
qui correspond au produit, \deleted{pour que ce soit}\inserted{il est bien}
\deleted{défini, il faudrait que pour tout $V\in K^{\underline U}$,}
\deleted{existe un $W\in K^{\underline U}$ tel que $\mathcal J_W U\subset$}
défini par passage au quotient).

\uline{Par dualité}, ceci correspond aux
données suivantes sur un ensemble $A$
\[
\left\{\vcenter{\hbox{\begin{minipage}{0.85\linewidth}
\deleted{\unsure{$A$}}: exactement les $=$s, sauf qu'il n'y a plus de topologie, et que c'est pour la structure d'algèbre ass. et non de co-algèbre que $A$ est \unsure{commutative|cocommutative}
\end{minipage}}}\right.
\]
% archive sheet number: 167
